% ═══════════════════════════════════════════════════════════════════════════
% Chapter 14: Error Analysis and Uncertainty Quantification
% ═══════════════════════════════════════════════════════════════════════════

\chapter{Error Analysis and Uncertainty Quantification} \label{chap:error_analysis}

\section{Overview}

This chapter provides a comprehensive error budget analysis, quantifying the contribution of each system component to the overall measurement uncertainty. Understanding error sources is essential for interpreting validation results and guiding future improvements.

\section{Error Sources Classification}

\subsection{Systematic Errors}

Errors that consistently bias measurements in one direction:
\begin{itemize}
    \item Accelerometer scale factor error
    \item Gyro bias drift over temperature
    \item Gravity model inaccuracy (local g variation)
    \item Filter steady-state offset
\end{itemize}

\subsection{Random Errors}

Errors with zero mean but non-zero variance:
\begin{itemize}
    \item Sensor white noise
    \item Quantization error
    \item Numerical round-off
    \item Timestamp jitter
\end{itemize}

\subsection{Environmental Errors}

Errors induced by operating conditions:
\begin{itemize}
    \item Temperature-dependent bias shift
    \item Vibration-induced measurement errors
    \item Electromagnetic interference
\end{itemize}

\section{Accelerometer Error Budget}

\subsection{ICM-42688-P Specifications}

\begin{table}[h]
    \centering
    \caption{Accelerometer error sources}
    \label{tab:accel_error}
    \begin{tabular}{lcc}
        \toprule
        \textbf{Source} & \textbf{Magnitude} & \textbf{Effect} \\
        \midrule
        Zero-g bias & $\pm$50 mg & Position drift \\
        Bias temperature coeff. & 0.5 mg/$^\circ$C & Slow drift \\
        Scale factor error & $\pm$1\% & ROM scaling error \\
        Scale factor temp. coeff. & 0.01\%/$^\circ$C & Minor \\
        Non-linearity & $\pm$1\% FS & High-accel distortion \\
        Noise density & 20 $\mu$g/$\sqrt{\text{Hz}}$ & High-frequency noise \\
        \bottomrule
    \end{tabular}
\end{table}

\subsection{Bias-Induced Position Error}

A constant bias $b_a$ produces quadratic position error:

\begin{equation}
    \epsilon_p(t) = \frac{1}{2} b_a t^2
\end{equation}

For $b_a = 5$ mg = 0.049 m/s\textsuperscript{2}:
\begin{itemize}
    \item After 1 s: $\epsilon_p = 2.5$ cm
    \item After 2 s: $\epsilon_p = 9.8$ cm
    \item After 5 s: $\epsilon_p = 61$ cm
\end{itemize}

This demonstrates the necessity of ZUPT for bounding drift.

\section{Gyroscope Error Budget}

\begin{table}[h]
    \centering
    \caption{Gyroscope error sources}
    \label{tab:gyro_error}
    \begin{tabular}{lcc}
        \toprule
        \textbf{Source} & \textbf{Magnitude} & \textbf{Effect} \\
        \midrule
        Zero-rate offset & $\pm$30 dps & Attitude drift \\
        Offset temp. coeff. & 0.1 dps/$^\circ$C & Slow drift \\
        Scale factor error & $\pm$1\% & Rotation rate error \\
        Non-linearity & $\pm$1\% FS & High-rate distortion \\
        Noise density & 0.015 deg/s/$\sqrt{\text{Hz}}$ & High-freq jitter \\
        \bottomrule
    \end{tabular}
\end{table}

\section{Orientation Filter Error Propagation}

\subsection{Quaternion Normalization Error}

Unit quaternion maintenance requires periodic normalization. Floating-point round-off can cause drift:

\begin{equation}
    \delta q = |1 - |q|| \approx \frac{1}{2}\sum \epsilon_i^2
\end{equation}

With double-precision arithmetic ($\epsilon \approx 10^{-16}$), this error is negligible.

\subsection{Accel Correction Weight Error}

During dynamic motion, accelerometer contains non-gravitational components. The VQF adaptive weight $w_a$ down-weights accelerometer trust. Residual error is:

\begin{equation}
    \epsilon_q \approx (1 - w_a) \cdot \theta_{motion}
\end{equation}

where $\theta_{motion}$ is the tilt caused by linear acceleration. For typical lifts with $\theta_{motion} \approx 15^\circ$ and $w_a \approx 0.2$, $\epsilon_q \approx 3^\circ$.

\section{Integration Numerical Error}

\subsection{Trapezoidal Rule Error}

Velocity integration uses trapezoidal rule:

\begin{equation}
    v_{k+1} = v_k + \frac{a_k + a_{k+1}}{2} \Delta t
\end{equation}

Local truncation error:

\begin{equation}
    \epsilon_v = \frac{1}{12} \dddot{v} (\Delta t)^3
\end{equation}

For $\Delta t = 1$ ms and $\dddot{v} \approx 100$ m/s\textsuperscript{3} (aggressive lift):

\begin{equation}
    \epsilon_v \approx 8 \times 10^{-6} \text{ m/s per sample}
\end{equation}

Over 1000 samples (1 s): accumulated error $\approx 8$ mm/s, negligible.

\section{Timestamp Uncertainty}

\subsection{SPI Read Latency}

SPI transactions complete in $45 \pm 3$ $\mu$s. This jitter introduces timestamp uncertainty negligible relative to 1 ms sample period.

\subsection{ESP-NOW Transmission Delay}

ESP-NOW latency is $3.2 \pm 0.3$ ms. The host records arrival time, but the IMU sample includes its own timestamp. End-to-end skew is dominated by:

\begin{equation}
    \epsilon_{timestamp} = \epsilon_{IMU\_clock} + \epsilon_{camera\_clock}
\end{equation}

With crystal accuracies of 50 ppm, over a 10 s session:

\begin{equation}
    \epsilon_{skew} = 50 \times 10^{-6} \times 10 = 0.5 \text{ ms}
\end{equation}

Synchronization calibration reduces this to $\approx 200$ $\mu$s.

\section{Ground Truth Uncertainty}

\subsection{Camera Tracking Error}

The RealSense D455 provides optical ground truth with stated depth accuracy:

\begin{itemize}
    \item At 1 m: $\pm$2 mm
    \item At 2 m: $\pm$5 mm
    \item At 3 m: $\pm$10 mm
\end{itemize}

For our setup (camera at $\approx$3 m), positional uncertainty is $\pm$10 mm.

\subsection{Velocity Computation Error}

Central difference velocity from discrete positions:

\begin{equation}
    v_k = \frac{p_{k+1} - p_{k-1}}{2\Delta t}
\end{equation}

With $\Delta t = 11.1$ ms (90 fps) and $\delta p = 10$ mm:

\begin{equation}
    \delta v = \frac{\delta p}{\Delta t} = \frac{0.01}{0.0111} \approx 0.9 \text{ m/s}
\end{equation}

This is unacceptably high; hence low-pass filtering (6 Hz cutoff) reduces noise at the cost of some bandwidth.

\section{Total Uncertainty Budget}

\begin{table}[h]
    \centering
    \caption{Combined uncertainty for peak velocity}
    \label{tab:total_uncertainty_vel}
    \begin{tabular}{lrr}
        \toprule
        \textbf{Source} & \textbf{Uncertainty} & \textbf{Combination} \\
        \midrule
        Accelerometer bias & 0.025 m/s & \multirow{4}{*}{RSS} \\
        Gyro drift & 0.015 m/s & \\
        Filter residual & 0.030 m/s & \\
        Integration numerical & 0.008 m/s & \\
        Timestamp skew & 0.005 m/s & \\
        \midrule
        \textbf{Combined (IMU)} & \textbf{0.041 m/s} & \\
        Ground truth uncertainty & 0.050 m/s & \\
        \midrule
        \textbf{Total uncertainty} & \textbf{0.065 m/s} & RSS \\
        \bottomrule
    \end{tabular}
\end{table}

Comparing to measured MAE of 0.068 m/s: error budget accounts for $\approx$95\% of observed error.

\begin{table}[h]
    \centering
    \caption{Combined uncertainty for range of motion}
    \label{tab:total_uncertainty_rom}
    \begin{tabular}{lrr}
        \toprule
        \textbf{Source} & \textbf{Uncertainty} & \textbf{Combination} \\
        \midrule
        Accelerometer scale & 15 mm & \multirow{5}{*}{RSS} \\
        Bias-induced drift & 20 mm & \\
        Orientation error & 10 mm & \\
        Integration numerical & 3 mm & \\
        Boundary timing & 8 mm & \\
        \midrule
        \textbf{Combined (IMU)} & \textbf{28 mm} & \\
        Ground truth uncertainty & 10 mm & \\
        \midrule
        \textbf{Total uncertainty} & \textbf{30 mm} & RSS \\
        \bottomrule
    \end{tabular}
\end{table}

Measured MAE of 23 mm is within the uncertainty budget.

\section{Sensitivity Analysis}

\subsection{Filter Parameter Sensitivity}

VQF parameters were tuned empirically. Sensitivity analysis shows:

\begin{table}[h]
    \centering
    \caption{VQF parameter sensitivity}
    \label{tab:filter_sensitivity}
    \begin{tabular}{lcc}
        \toprule
        \textbf{Parameter} & \textbf{Variation} & \textbf{MAE change} \\
        \midrule
        $\tau_{acc}$ & $\pm$50\% & $\pm$0.012 m/s \\
        $\tau_{bias}$ & $\pm$50\% & $\pm$0.008 m/s \\
        Rest accel gate & $\pm$20\% & $\pm$0.015 m/s \\
        Rest gyro gate & $\pm$20\% & $\pm$0.010 m/s \\
        \bottomrule
    \end{tabular}
\end{table}

Algorithm is moderately sensitive to tuning; VQF defaults work well across exercises.

\section{Outlier Analysis}

\subsection{Identifying Outliers}

Reps with error $>$ 3 standard deviations are flagged:

\begin{itemize}
    \item 12 of 622 reps (1.9\%) classified as outliers
    \item Primary causes: partial reps, dropped samples, video occlusion
\end{itemize}

\subsection{Excluding Outliers}

With outliers removed:
\begin{itemize}
    \item Peak velocity MAE: 0.058 m/s (vs. 0.068 m/s)
    \item ROM MAE: 19 mm (vs. 23 mm)
\end{itemize}

\section{Monte Carlo Simulation}

\subsection{Experimental Setup}

10,000 synthetic rep traces generated with:
\begin{itemize}
    \item Random velocity profile within physiological bounds
    \item Additive sensor noise per specifications
    \item Random bias offset per temperature range
\end{itemize}

\subsection{Results}

Simulation predicts:
\begin{itemize}
    \item Velocity MAE: 0.062 $\pm$ 0.005 m/s
    \item ROM MAE: 25 $\pm$ 3 mm
\end{itemize}

Consistent with empirical measurements.

\section{Summary}

Error analysis confirms that:
\begin{itemize}
    \item IMU sensor errors dominate uncertainty budget
    \item Orientation filter residuals are well-bounded by VQF
    \item Ground truth contributes $\approx$20\% of total uncertainty
    \item Outliers (2\% of data) skew aggregate metrics slightly
\end{itemize}

The measured performance (0.068 m/s velocity MAE, 23 mm ROM MAE) aligns with theoretical error budget predictions, validating the system design.

\newpage